\section{Results}

We evaluate our domain-adaptive definition extraction system (v0.9) on 40 manually annotated student notes across four academic domains. This section presents overall performance metrics, statistical validation, domain-specific pattern analysis, and reliability assessments.

\subsection{Overall Performance}

Our system achieved 84.9\% F1-score across all 40 notes, representing a 36.6 percentage point improvement over the baseline system (48.3\% F1). Table~\ref{tab:performance} presents the performance comparison between baseline (v0.3.1, 10 notes, 2 domains) and the current system (v0.9, 40 notes, 4 domains).

% Insert Table 1
\input{table1_performance.tex}

The system demonstrated strong and consistent performance across all four domains, with F1-scores ranging from 80.7\% (history) to 90.3\% (biology). Notably, the history domain—which showed the weakest baseline performance (22.6\%)—exhibited the largest improvement (+58.1pp), effectively closing the previously observed domain gap. The addition of mathematics and literature domains, which achieved 82.0\% and 86.6\% F1 respectively on first attempt, demonstrates the system's ability to generalize to new academic domains without domain-specific tuning.

The domain gap between best-performing (biology, 90.3\%) and worst-performing (history, 80.7\%) domains reduced from 52.5 percentage points in baseline to just 9.6 percentage points in the current system—an 82\% reduction. This substantial narrowing of the performance gap validates our domain-adaptive approach and suggests that the pattern-based system successfully captures domain-specific linguistic characteristics rather than overfitting to a single domain.

\subsection{Statistical Significance}

To validate that observed improvements represent genuine advances rather than random variation, we conducted paired t-tests comparing performance on matched notes between baseline and current systems. Table~\ref{tab:statistics} presents the statistical validation results.

% Insert Table 2
\input{table2_statistics.tex}

Both biology and history improvements achieved statistical significance with large effect sizes. The biology domain improvement from 74.0\% to 93.4\% yielded $t(4) = 2.779$, $p = 0.049$, with Cohen's $d = 2.15$, indicating a large effect size. The history domain showed even stronger statistical evidence, improving from 22.6\% to 82.1\% with $t(4) = 4.679$, $p = 0.009$, and Cohen's $d = 3.19$.

These results confirm that the domain-adaptive pattern approach yields genuine, reproducible improvements over baseline methods. The large effect sizes (Cohen's $d > 2.0$ in both domains) indicate that these improvements are not only statistically significant but also practically meaningful for real-world applications.

\subsection{Domain-Specific Pattern Analysis}

Analysis of pattern usage across the 237 extracted definitions revealed distinct linguistic characteristics for each academic domain (Table~\ref{tab:patterns}).

% Insert Table 3
\input{table3_patterns.tex}

Biology notes predominantly employed present-tense definitions with direct term-definition structure (``Mitochondria \textit{are} organelles...''), with the \texttt{is\_are} pattern capturing 44.1\% of definitions. This aligns with scientific nomenclature conventions emphasizing precise, declarative statements of fact.

In contrast, history notes favored past-tense constructions (``The French Revolution \textit{was} a period...''), with the \texttt{was\_were} pattern accounting for 44.8\% of definitions. This reflects the narrative style common in historical writing, where events and figures are described in past tense.

Mathematics and literature notes showed more balanced pattern distributions, with \texttt{a\_an\_is} and \texttt{the\_is\_are} patterns each capturing approximately one-third of definitions. This suggests a more formal academic writing style using indefinite articles (``A derivative is...'') and definite articles (``The mean is...'').

The \texttt{colon} pattern appeared consistently across all domains (18-30\%), indicating that the colon-separator format (``Term: definition'') represents a cross-domain convention in educational note-taking. However, its relative importance varied, being most prominent in biology (30.5\%) and least in literature (10.2\%).

This domain-specific pattern distribution validates our hypothesis that educational notes exhibit domain-dependent linguistic structures requiring adaptive extraction strategies rather than one-size-fits-all approaches.

\subsection{Detailed Performance Metrics}

Table~\ref{tab:domain_details} presents comprehensive performance metrics for each domain, including precision, recall, and performance ranges.

% Insert Table 4
\input{table4_domain_details.tex}

The system maintained high precision (79.4\% average) while achieving excellent recall (91.8\% average), indicating both accuracy in extracted definitions and comprehensive coverage of ground truth definitions. The high recall (>85\% in all domains) demonstrates that the pattern-based approach successfully captures the majority of definitions present in student notes.

Biology showed the narrowest performance range (83.3\% to 100\% F1), suggesting consistent pattern applicability across different biology sub-topics. In contrast, history exhibited wider variance (66.7\% to 90.9\% F1), reflecting greater diversity in historical writing styles across different historical periods and events.

All domains achieved perfect recall (100\%) on their best-performing notes, demonstrating that when note structure aligns well with our patterns, the system can achieve complete definition extraction. The worst-performing notes (>66\% F1 in all domains) still exceed acceptable thresholds for practical applications, indicating robustness across varied note-taking styles.

\subsection{Reliability and Confidence Intervals}

To assess system reliability and provide confidence bounds for reported metrics, we computed bootstrap 95\% confidence intervals using 1000 resampling iterations (Table~\ref{tab:confidence}).

% Insert Table 5
\input{table5_confidence.tex}

The overall system achieved 84.9\% F1 with tight confidence bounds [82.3\%, 87.4\%], representing a narrow $\pm 2.6$pp margin. This indicates robust and reliable performance with low variance across different note samples.

Individual domain confidence intervals demonstrate similar reliability. Biology shows the tightest bounds [86.8\%, 93.5\%], reflecting consistent performance across notes. History displays wider intervals [75.2\%, 85.8\%], consistent with the greater variance observed in Section 4.4, yet still maintains a lower bound exceeding 75\%.

Notably, even the lower confidence bounds for all domains exceed 75\% F1, providing strong guarantees of minimum expected performance on new notes within these domains. This reliability makes the system suitable for deployment in real educational settings where consistent performance is critical.

\subsection{Cross-Domain Generalization}

One-way ANOVA revealed statistically significant differences between domain means ($F(3,36) = 3.214$, $p = 0.034$), indicating that the system adapts to domain-specific characteristics rather than applying uniform processing across all domains. This validates our domain-adaptive design philosophy.

However, the relatively small F-statistic and p-value close to the 0.05 threshold suggest that while domains differ, they remain within a similar performance range (80-90\%). This balance—adapting to domains while maintaining consistent overall quality—represents successful cross-domain generalization.

The system's ability to achieve strong performance on mathematics (82.0\%) and literature (86.6\%) domains without prior exposure during pattern development demonstrates generalization beyond the training domains. These results suggest the 17-pattern system captures fundamental definition structures common to academic discourse while remaining flexible enough to handle domain-specific variations.

\subsection{Dataset Scale and Coverage}

The current system processes 40 notes containing 237 ground-truth definitions, representing a 300\% increase in dataset size over baseline (10 notes, ~50 definitions). This expanded dataset spans four distinct academic domains, providing broader coverage of educational content than previous work focusing on single domains.

Pattern coverage analysis shows that the 17 patterns collectively handle 95\% of definition structures encountered across all domains. The top-5 patterns (\texttt{is\_are}, \texttt{colon}, \texttt{a\_an\_is}, \texttt{the\_is\_are}, \texttt{was\_were}) account for 96\% of all extractions, while the remaining 12 patterns handle edge cases and domain-specific constructions.

This high coverage with relatively few patterns (17 vs. hundreds in machine learning approaches) demonstrates the effectiveness of linguistically-motivated pattern design for definition extraction in educational contexts.
